\documentclass{article}
\usepackage[margin=1.5cm]{geometry}
\usepackage{amsmath}
\usepackage{enumitem}

\setlist{nosep,leftmargin=*}

\begin{document}
\twocolumn

\title{Problem Set 3: Loops and Functions}
\author{Prof. Jordan C. Hanson}
\date{}
\maketitle

\section{Exercise 1: Evaluating a Function}

\begin{enumerate}
\item Define the mathematical function
\[
f(x)=\frac{x+1}{x-1}.
\]
The function has a pole at $x=1$.  Write a Python function

\begin{verbatim}
def f(x):
    return (x + 1) / (x - 1)
\end{verbatim}

Use the list

\begin{verbatim}
values = [0, 0.5, 0.9,
          1.1, 1.5, 2]
\end{verbatim}

and a \texttt{for} loop to evaluate and print $x$ and $f(x)$ for every value in the list.

\begin{itemize}
\item Your loop must call \texttt{f(x)} rather than repeating the formula inside the loop.
\item Explain why $x=1$ is not included in \texttt{values}.
\item Add the values $0.99$ and $1.01$.  What happens to the magnitude and sign of $f(x)$ near the pole?
\end{itemize}
\end{enumerate}

\section{Exercise 2: Placement Function}

\begin{enumerate}
\item A placement score determines the appropriate math course:
\begin{center}
\begin{tabular}{c|c}
\textbf{Score} & \textbf{Course} \\ \hline
0--1 & Algebra \\
2--3 & Trigonometry \\
4--5 & Calculus
\end{tabular}
\end{center}

Write a function

\begin{verbatim}
def placement(score):
    # use if, elif, else
    # return a course name
\end{verbatim}

that \textbf{returns} the appropriate course name as a string.
Test the function first with

\begin{verbatim}
print(placement(1))
print(placement(3))
print(placement(5))
\end{verbatim}

The results should be Algebra, Trigonometry, and Calculus.

Now use

\begin{verbatim}
students = [
 ("Maya", 20681001, 5),
 ("Noah", 20681002, 1),
 ("Olivia", 20681003, 3),
 ("Liam", 20681004, 4),
 ("Sophia", 20681005, 2),
 ("Ethan", 20681006, 0)
]
\end{verbatim}

Create the roster structure

\begin{verbatim}
algebra = ["Algebra", []]
trig = ["Trigonometry", []]
calculus = ["Calculus", []]
math_courses = (algebra, trig, calculus)
\end{verbatim}

Use one \texttt{for} loop to process every student.  For each student:

\begin{itemize}
\item extract the name, student ID, and score;
\item call \texttt{placement(score)};
\item append the tuple \texttt{(name, student\_id)} to the correct roster.
\end{itemize}

Do not repeat the placement rules inside the loop; the decision must be made by the function.

Finally, print \texttt{math\_courses} and verify that every student appears exactly once and in the correct course.
\end{enumerate}

\section{Exercise 3: While-Loop Count}

\begin{enumerate}
\item Use a simple \texttt{while} loop to count the total number of students stored in the three rosters.

Begin with

\begin{verbatim}
i = 0
total = 0
\end{verbatim}

and use \texttt{i} to access each element of \texttt{math\_courses}.  During each pass, add the length of that course's student list to \texttt{total}, then increase \texttt{i} by one.

Print \texttt{total}.  What value should the program produce, and why?
\end{enumerate}

\end{document}
